\section*{Abstract}

This thesis asks to what extent current market and regulatory conditions produce a divergence between the commercial and socio-economic value of thermal flexibility in Norwegian district heating, and what mechanisms could reduce that divergence. Oslo's district heating system, operated by Hafslund Celsio, is the single case, with electric boilers, heat pumps and accumulator thermal storage as the flexibility assets.

The commercial side is measured with a deterministic linear dispatch model of the operator's own fleet, solved hour by hour over a delivery year as a cost minimiser under a fixed heat-delivery obligation and as a price taker. The model is run across two grid-tariff structures, the interruptible arrangement the operator holds and the standard schedule with a monthly capacity charge, crossed with three levels of reserve-market access, the 29 MW held today, 100 MW and the 200 MW applied for, on the 2023, 2024 and 2025 delivery years. The socio-economic side is triangulated from published unit values applied to the physical quantities the model produces, on two of six channels: freed grid capacity in the tightest one per cent of hours ranked by heat demand, and displaced fossil peaking net of the carbon levy. Because the socio-economic total contains the operator's own earnings, the divergence is measured as the division of that total between what the operator captures and what reaches other parties through no mechanism. A second quantity, the value that the current allocation of reserve-market access stops from being created at all, is reported beside that division, and the two are never summed.

Against the plant as it runs, flexibility is worth 65.79 million Norwegian kroner (MNOK) per year to the operator at the 2024 baseline, and 140.87 MNOK against a price-blind plant; the difference is price-responsive dispatch, which sits outside the headline because the plant already dispatches on price. On the two priced channels the operator captures 90 per cent of a priced socio-economic total of 73.05 MNOK per year, and the un-captured externality is 7.26 MNOK per year, between 5.41 and 10.95 across the band of published grid-valuation methods. That total is a floor, because the three unpriced channels are positive by construction. Beside it, the access allocation withholds 138.42 MNOK per year of upward reserve revenue between 29 MW and 200 MW, roughly nineteen times the externality; it is a transfer from the reserve market, and an upper bound because price taking fails at that volume. The interruptible arrangement is worth 45.68 MNOK per year at today's access and 84.90 at 200 MW, so the tariff and access barriers compound, and more than half of what the standard tariff costs the operator is the dispatch the charge provokes. Every sign and ordering holds in all three delivery years while the magnitudes move by factors between 1.7 and 2.9. What drives the variation is whether cheap power arrived while heat was wanted; in 2023 it did not, and the externality is negative, at $-$7.85 MNOK per year.

The divergence in Oslo is therefore mainly a blockage upstream of any transfer. The instruments that would have to move form a ladder from statute at the top to the system operator's own prequalification practice at the bottom, and the measured magnitudes run the other way along it. The capacity-area cap, lowest on the ladder, carries the largest quantity in the thesis, while the 22 to 33 MW band the current regime argues over sits on the flat top of the access curve. Because the two barriers interact, removing the cap without touching the tariff would hand back less than the access figures alone suggest. This is a ranking, and no recommendation is made.

The evidence base is one case, one operator and one interview. What the operator captures is overstated by at least a quarter of the baseline flexibility value through foresight of reserve prices and activation, an overstatement that falls on the operator's own value and leaves the ordering of the two registers intact. The sign of the un-captured externality turns on how scarce hours are defined, and is positive in twelve of the ninety combinations of definition and run that were swept, so every conclusion about it carries the definition that produced it.
